\documentclass[12pt,a4paper]{article}
\usepackage{amsmath,amssymb,amsthm}
\usepackage{geometry}
\usepackage{enumitem}
\usepackage{ctex}

\geometry{a4paper, margin=2.5cm}

% 微分符号
\newcommand{\dd}{\,\mathrm{d}}
% 极限符号
\newcommand{\Lim}[2]{\lim\limits_{#1 \to #2}}

\title{\textbf{数学分析习题解答}}
\author{}
\date{}

\begin{document}
\maketitle
\thispagestyle{empty}

\section*{习题一 极限与积分}

\begin{enumerate}[label=\textbf{\arabic*.}, leftmargin=*, itemsep=1.2em]

% 第1题
\item 求极限 $\displaystyle \Lim{n}{\infty} \left( \frac{1}{n^2} + \frac{2}{n^2} + \dots + \frac{n}{n^2} \right)$

\textbf{解：}
\begin{align*}
\text{原式} &= \Lim{n}{\infty} \frac{1+2+\dots+n}{n^2} \\
&= \Lim{n}{\infty} \frac{\frac{n(n+1)}{2}}{n^2} \\
&= \Lim{n}{\infty} \frac{n+1}{2n} = \frac{1}{2}
\end{align*}

% 第2题
\item 求极限 $\displaystyle \Lim{n}{\infty} \left( \frac{1}{n^3} + \frac{2^2}{n^3} + \dots + \frac{n^2}{n^3} \right)$

\textbf{解：}
\begin{align*}
\text{原式} &= \Lim{n}{\infty} \frac{1^2+2^2+\dots+n^2}{n^3} \\
&= \Lim{n}{\infty} \frac{\frac{n(n+1)(2n+1)}{6}}{n^3} \\
&= \Lim{n}{\infty} \frac{(n+1)(2n+1)}{6n^2} = \frac{1}{3}
\end{align*}

% 第3题
\item 求极限 $\displaystyle \Lim{n}{\infty} \left( \frac{1}{\sqrt{n^2+1}} + \frac{1}{\sqrt{n^2+2}} + \dots + \frac{1}{\sqrt{n^2+n}} \right)$

\textbf{解：}
\begin{align*}
\text{放缩：} \quad & \frac{n}{\sqrt{n^2+n}} \leq \text{原式} \leq \frac{n}{\sqrt{n^2+1}} \\
\Lim{n}{\infty} \frac{n}{\sqrt{n^2+n}} &= \Lim{n}{\infty} \frac{1}{\sqrt{1+\frac{1}{n}}} = 1 \\
\Lim{n}{\infty} \frac{n}{\sqrt{n^2+1}} &= \Lim{n}{\infty} \frac{1}{\sqrt{1+\frac{1}{n^2}}} = 1 \\
\text{由夹逼准则，原式} &= 1
\end{align*}

% 第4题
\item 求极限 $\displaystyle \Lim{n}{\infty} \left( 1 + \frac{1}{1+1} \right) \left( 1 + \frac{1}{2+1} \right) \dots \left( 1 + \frac{1}{n+1} \right)$

\textbf{解：}
\begin{align*}
\text{通项化简：} \quad 1 + \frac{1}{k+1} &= \frac{k+2}{k+1} \\
\text{乘积展开：} \quad & \frac{3}{2} \cdot \frac{4}{3} \cdot \frac{5}{4} \dots \frac{n+2}{n+1} = \frac{n+2}{2} \\
\text{原式} &= \Lim{n}{\infty} \frac{n+2}{2} = +\infty
\end{align*}

% 第5题
\item 设 $f(x)$ 在 $[a,b]$ 上可积，且 $\displaystyle \int_a^b f(x)\dd x = 1$，求 $\displaystyle \Lim{n}{\infty} \frac{1}{n} \sum_{k=1}^n f\left( \frac{k}{n} \right)$

\textbf{解：}
\begin{align*}
\text{由定积分定义：} \quad & \int_0^1 f(x)\dd x = \Lim{n}{\infty} \frac{1}{n} \sum_{k=1}^n f\left( \frac{k}{n} \right) \\
\text{原式} &= \int_0^1 f(x)\dd x = 1
\end{align*}

% 第6题
\item 求极限 $\displaystyle \Lim{n}{\infty} \sum_{k=1}^n \frac{k}{n^2 + k^2}$

\textbf{解：}
\begin{align*}
\sum_{k=1}^n \frac{k}{n^2 + k^2} &= \frac{1}{n} \sum_{k=1}^n \frac{\frac{k}{n}}{1 + \left( \frac{k}{n} \right)^2} \\
\text{原式} &= \int_0^1 \frac{x}{1+x^2} \dd x = \frac{1}{2} \ln(1+x^2) \bigg|_0^1 = \frac{1}{2}\ln 2
\end{align*}

% 第7题
\item 设 $\displaystyle I_n = \int_0^{\pi/2} \sin^n x \, \dd x$，证明 $I_n = \frac{n-1}{n} I_{n-2}$ 并求极限

\textbf{解：}
\begin{align*}
\text{分部积分：} \quad I_n &= -\cos x \sin^{n-1}x \bigg|_0^{\pi/2} + (n-1) \int_0^{\pi/2} \sin^{n-2}x \cos^2x \, \dd x \\
&= (n-1) \int_0^{\pi/2} \sin^{n-2}x (1-\sin^2x) \, \dd x \\
&= (n-1)I_{n-2} - (n-1)I_n \\
\Rightarrow nI_n &= (n-1)I_{n-2} \Rightarrow I_n = \frac{n-1}{n}I_{n-2}
\end{align*}

\textbf{极限：} $\displaystyle \Lim{n}{\infty} I_n = 0$（因 $\sin^n x$ 一致收敛于0）

% 第8题
\item 求 $\displaystyle \int_0^{\pi/2} \sin^n x \, \dd x$ 的递推公式与具体值

\textbf{解：}
\begin{align*}
n \text{ 偶：} \quad I_n &= \frac{n-1}{n} \cdot \frac{n-3}{n-2} \dots \frac{1}{2} \cdot \frac{\pi}{2} \\
n \text{ 奇：} \quad I_n &= \frac{n-1}{n} \cdot \frac{n-3}{n-2} \dots \frac{2}{3} \cdot 1
\end{align*}

% 第9题
\item 求 $\displaystyle \int x e^x \, \dd x$

\textbf{解：}
\begin{align*}
u=x, \dd v=e^x \dd x &\Rightarrow \dd u=\dd x, v=e^x \\
\int x e^x \dd x &= x e^x - \int e^x \dd x = x e^x - e^x + C
\end{align*}

% 第10题
\item 求 $\displaystyle \int x \cos x \, \dd x$

\textbf{解：}
\begin{align*}
u=x, \dd v=\cos x \dd x &\Rightarrow \dd u=\dd x, v=\sin x \\
\int x \cos x \dd x &= x \sin x - \int \sin x \dd x = x \sin x + \cos x + C
\end{align*}

% 第11题
\item 求 $\displaystyle \int x^2 e^x \, \dd x$

\textbf{解：}
\begin{align*}
u=x^2, \dd v=e^x \dd x &\Rightarrow \dd u=2x \dd x, v=e^x \\
\int x^2 e^x \dd x &= x^2 e^x - 2 \int x e^x \dd x \\
&= x^2 e^x - 2(x e^x - e^x) + C = e^x(x^2-2x+2)+C
\end{align*}

% 第12题
\item 求 $\displaystyle \int \ln x \, \dd x$

\textbf{解：}
\begin{align*}
u=\ln x, \dd v=\dd x &\Rightarrow \dd u=\frac{1}{x}\dd x, v=x \\
\int \ln x \dd x &= x \ln x - \int 1 \, \dd x = x \ln x - x + C
\end{align*}

% 第13题
\item 求 $\displaystyle \int x \ln x \, \dd x$

\textbf{解：}
\begin{align*}
u=\ln x, \dd v=x \dd x &\Rightarrow \dd u=\frac{1}{x}\dd x, v=\frac{x^2}{2} \\
\int x \ln x \dd x &= \frac{x^2}{2}\ln x - \int \frac{x^2}{2} \cdot \frac{1}{x} \dd x = \frac{x^2}{2}\ln x - \frac{x^2}{4} + C
\end{align*}

% 第14题
\item 求 $\displaystyle \int \sin^2 x \, \dd x$

\textbf{解：}
\begin{align*}
\sin^2 x &= \frac{1-\cos 2x}{2} \\
\int \sin^2 x \dd x &= \frac{1}{2} \int (1-\cos 2x) \dd x = \frac{x}{2} - \frac{\sin 2x}{4} + C
\end{align*}

% 第15题
\item 求 $\displaystyle \int \cos^2 x \, \dd x$

\textbf{解：}
\begin{align*}
\cos^2 x &= \frac{1+\cos 2x}{2} \\
\int \cos^2 x \dd x &= \frac{1}{2} \int (1+\cos 2x) \dd x = \frac{x}{2} + \frac{\sin 2x}{4} + C
\end{align*}

% 第16题
\item 求 $\displaystyle \int_0^{\pi} \sin^3 x \, \dd x$

\textbf{解：}
\begin{align*}
\sin^3 x &= \sin x (1-\cos^2 x) \\
\int_0^{\pi} \sin^3 x \dd x &= -\int_0^{\pi} (1-\cos^2 x) \dd(\cos x) \\
&= -\left( \cos x - \frac{\cos^3 x}{3} \right) \bigg|_0^{\pi} = -\left( (-1 + \frac{1}{3}) - (1 - \frac{1}{3}) \right) = \frac{4}{3}
\end{align*}

% 第17题
\item 求 $\displaystyle \int \frac{1}{x^2 + a^2} \dd x$ $(a>0)$

\textbf{解：}
\begin{align*}
\int \frac{1}{x^2 + a^2} \dd x &= \frac{1}{a^2} \int \frac{1}{1 + \left( \frac{x}{a} \right)^2} \dd x = \frac{1}{a} \arctan\left( \frac{x}{a} \right) + C
\end{align*}

% 第18题
\item 求 $\displaystyle \int \frac{1}{\sqrt{a^2 - x^2}} \dd x$ $(a>0)$

\textbf{解：}
\begin{align*}
\int \frac{1}{\sqrt{a^2 - x^2}} \dd x &= \arcsin\left( \frac{x}{a} \right) + C
\end{align*}

% 第19题
\item 求 $\displaystyle \int \frac{1}{x^2 - a^2} \dd x$ $(a>0)$

\textbf{解：}
\begin{align*}
\frac{1}{x^2 - a^2} &= \frac{1}{2a} \left( \frac{1}{x-a} - \frac{1}{x+a} \right) \\
\int \frac{1}{x^2 - a^2} \dd x &= \frac{1}{2a} \ln \left| \frac{x-a}{x+a} \right| + C
\end{align*}

% 第20题
\item 求 $\displaystyle \int \sec x \, \dd x$

\textbf{解：}
\begin{align*}
\int \sec x \dd x &= \ln |\sec x + \tan x| + C
\end{align*}

% 第21题
\item 求 $\displaystyle \int \csc x \, \dd x$

\textbf{解：}
\begin{align*}
\int \csc x \dd x &= \ln |\csc x - \cot x| + C
\end{align*}

% 第22题
\item 求 $\displaystyle \int_0^{\pi/2} \frac{\sin x}{\sin x + \cos x} \dd x$

\textbf{解：}
\begin{align*}
I &= \int_0^{\pi/2} \frac{\sin x}{\sin x + \cos x} \dd x = \int_0^{\pi/2} \frac{\cos x}{\cos x + \sin x} \dd x \quad (\text{换元 } x=\pi/2-t) \\
2I &= \int_0^{\pi/2} 1 \, \dd x = \frac{\pi}{2} \Rightarrow I = \frac{\pi}{4}
\end{align*}

% 第23题
\item 求 $\displaystyle \int \frac{x}{x^2 + 2x + 2} \dd x$

\textbf{解：}
\begin{align*}
\text{配方：} x^2+2x+2 &= (x+1)^2 + 1 \\
\int \frac{x}{(x+1)^2 + 1} \dd x &= \frac{1}{2} \ln(x^2+2x+2) - \arctan(x+1) + C
\end{align*}

% 第24题
\item 求 $\displaystyle \int \sqrt{x^2 + a^2} \dd x$ $(a>0)$

\textbf{解：}
\begin{align*}
\text{分部积分：} \quad u=\sqrt{x^2+a^2}, \dd v=\dd x \\
\int \sqrt{x^2+a^2} \dd x &= \frac{x}{2}\sqrt{x^2+a^2} + \frac{a^2}{2} \ln\left( x + \sqrt{x^2+a^2} \right) + C
\end{align*}

% 第25题
\item 求 $\displaystyle \int \sqrt{a^2 - x^2} \dd x$ $(a>0)$

\textbf{解：}
\begin{align*}
\text{三角换元：} x=a\sin t \\
\int \sqrt{a^2 - x^2} \dd x &= \frac{a^2}{2} \arcsin\frac{x}{a} + \frac{x}{2}\sqrt{a^2 - x^2} + C
\end{align*}

% 第26题
\item 求 $\displaystyle \int \frac{1}{\sqrt{x^2 + 2x}} \dd x$

\textbf{解：}
\begin{align*}
\text{配方：} x^2+2x &= (x+1)^2 - 1 \\
\int \frac{1}{\sqrt{(x+1)^2 - 1}} \dd x &= \ln \left| x+1 + \sqrt{x^2+2x} \right| + C
\end{align*}

% 第27题
\item 求 $\displaystyle \int_0^1 x e^{-x} \dd x$

\textbf{解：}
\begin{align*}
u=x,\ \dd v=e^{-x}\dd x &\Rightarrow \dd u=\dd x,\ v=-e^{-x}\\
\int_0^1 x e^{-x} \dd x &= -x e^{-x}\big|_0^1 + \int_0^1 e^{-x} \dd x\\
&= -e^{-1} - e^{-x}\big|_0^1 = -\frac{1}{e} - \left(\frac{1}{e}-1\right) = 1 - \frac{2}{e}
\end{align*}

% 第28题
\item 求 $\displaystyle \int \frac{\ln x}{x^2} \dd x$

\textbf{解：}
\begin{align*}
u=\ln x,\ \dd v&=\frac{1}{x^2}\dd x \Rightarrow \dd u=\frac{1}{x}\dd x,\ v=-\frac{1}{x}\\
\int \frac{\ln x}{x^2} \dd x &= -\frac{\ln x}{x} + \int \frac{1}{x^2} \dd x = -\frac{\ln x}{x} - \frac{1}{x} + C
\end{align*}

% 第29题
\item 求 $\displaystyle \int e^{\sqrt{x}} \dd x$

\textbf{解：}
\begin{align*}
\text{令 } t=\sqrt{x} &\Rightarrow x=t^2,\ \dd x=2t \dd t\\
\int e^{\sqrt{x}} \dd x &= 2\int t e^t \dd t = 2\left(t e^t - e^t\right) + C = 2e^{\sqrt{x}}(\sqrt{x}-1) + C
\end{align*}

% 第30题
\item 求 $\displaystyle \int_0^1 \arcsin x \dd x$

\textbf{解：}
\begin{align*}
u=\arcsin x,\ \dd v&=\dd x \Rightarrow \dd u=\frac{1}{\sqrt{1-x^2}}\dd x,\ v=x\\
\int_0^1 \arcsin x \dd x &= x\arcsin x\big|_0^1 - \int_0^1 \frac{x}{\sqrt{1-x^2}} \dd x\\
&= \frac{\pi}{2} + \sqrt{1-x^2}\big|_0^1 = \frac{\pi}{2} - 1
\end{align*}

% 第31题
\item 求 $\displaystyle \int \arctan x \dd x$

\textbf{解：}
\begin{align*}
u=\arctan x,\ \dd v&=\dd x \Rightarrow \dd u=\frac{1}{1+x^2}\dd x,\ v=x\\
\int \arctan x \dd x &= x\arctan x - \frac{1}{2}\int \frac{1}{1+x^2}\dd(1+x^2)\\
&= x\arctan x - \frac{1}{2}\ln(1+x^2) + C
\end{align*}

% 第32题
\item 设 $\displaystyle F(x)=\int_0^x te^{-t^2}\dd t$，求 $F'(x)$ 与极值

\textbf{解：}
\begin{align*}
F'(x) &= x e^{-x^2}\\
F''(x) &= e^{-x^2}(1-2x^2)\\
\text{令 } F'(x)=0 &\Rightarrow x=0\\
F''(0) &= 1>0 \Rightarrow x=0 \text{ 为极小值点，极小值 } F(0)=0
\end{align*}

% 第33题
\item 求 $\displaystyle \frac{\dd}{\dd x}\int_0^{x^2} \sqrt{1+t^2}\dd t$

\textbf{解：}
\begin{align*}
\text{由变上限求导：} \quad \frac{\dd}{\dd x}\int_0^{x^2} \sqrt{1+t^2}\dd t = \sqrt{1+x^4}\cdot 2x = 2x\sqrt{1+x^4}
\end{align*}

% 第34题
\item 求 $\displaystyle \frac{\dd}{\dd x}\int_x^{x^2} e^{-t^2}\dd t$

\textbf{解：}
\begin{align*}
\text{原式} &= e^{-x^4}\cdot 2x - e^{-x^2}\cdot 1 = 2x e^{-x^4} - e^{-x^2}
\end{align*}

% 第35题
\item 求极限 $\displaystyle \lim\limits_{x\to 0}\frac{\int_0^x \sin t^2 \dd t}{x^3}$

\textbf{解：}
\begin{align*}
\text{洛必达：} \quad \lim_{x\to 0}\frac{\sin x^2}{3x^2} = \lim_{x\to 0}\frac{x^2}{3x^2} = \frac{1}{3}
\end{align*}

% 第36题
\item 求极限 $\displaystyle \lim\limits_{x\to 0}\frac{\int_0^x \ln(1+t)\dd t}{x^2}$

\textbf{解：}
\begin{align*}
\text{洛必达：} \quad \lim_{x\to 0}\frac{\ln(1+x)}{2x} = \lim_{x\to 0}\frac{x}{2x} = \frac{1}{2}
\end{align*}

% 第37题
\item 求 $\displaystyle \int_0^{+\infty} e^{-x}\dd x$

\textbf{解：}
\begin{align*}
\int_0^{+\infty} e^{-x}\dd x &= \lim\limits_{b\to +\infty} \left(-e^{-x}\right)\big|_0^b = 0 - (-1) = 1
\end{align*}

% 第38题
\item 求 $\displaystyle \int_1^{+\infty} \frac{1}{x^2}\dd x$

\textbf{解：}
\begin{align*}
\int_1^{+\infty} \frac{1}{x^2}\dd x &= \lim\limits_{b\to +\infty} \left(-\frac{1}{x}\right)\big|_1^b = 0 - (-1) = 1
\end{align*}

% 第39题
\item 求 $\displaystyle \int_0^{+\infty} x e^{-x}\dd x$

\textbf{解：}
\begin{align*}
\int_0^{+\infty} x e^{-x}\dd x &= \lim\limits_{b\to +\infty} \left(-x e^{-x}-e^{-x}\right)\big|_0^b = 0 - (-1) = 1
\end{align*}

% 第40题
\item 求 $\displaystyle \int_0^1 \frac{1}{\sqrt{x}}\dd x$

\textbf{解：}
\begin{align*}
\int_0^1 \frac{1}{\sqrt{x}}\dd x &= \lim\limits_{\varepsilon\to 0^+} 2\sqrt{x}\big|_\varepsilon^1 = 2 - 0 = 2
\end{align*}

% 第41题
\item 求 $\displaystyle \int_0^1 \ln x \dd x$

\textbf{解：}
\begin{align*}
\int_0^1 \ln x \dd x &= \lim\limits_{\varepsilon\to 0^+} \left(x\ln x - x\right)\big|_\varepsilon^1 = -1 - 0 = -1
\end{align*}

% 第42题
\item 求由 $y=x^2$ 与 $y=x$ 围成的面积

\textbf{解：}
\begin{align*}
&\text{交点：}x=0, x=1\\
S &= \int_0^1 (x-x^2)\dd x = \left(\frac{x^2}{2}-\frac{x^3}{3}\right)\big|_0^1 = \frac{1}{2}-\frac{1}{3} = \frac{1}{6}
\end{align*}

% 第43题
\item 求由 $y=e^x, y=e^{-x}, x=1$ 围成的面积

\textbf{解：}
\begin{align*}
S &= \int_0^1 (e^x-e^{-x})\dd x = \left(e^x+e^{-x}\right)\big|_0^1 = e+\frac{1}{e}-2
\end{align*}

% 第44题
\item 求 $y=x^2$ 绕 $x$ 轴在 $[0,1]$ 上的旋转体体积

\textbf{解：}
\begin{align*}
V_x &= \pi\int_0^1 (x^2)^2 \dd x = \pi\int_0^1 x^4 \dd x = \frac{\pi}{5}
\end{align*}

% 第45题
\item 求 $y=\sin x\ (0\le x\le\pi)$ 绕 $x$ 轴旋转体体积

\textbf{解：}
\begin{align*}
V_x &= \pi\int_0^\pi \sin^2 x \dd x = \pi\int_0^\pi \frac{1-\cos 2x}{2}\dd x = \frac{\pi^2}{2}
\end{align*}

\end{enumerate}

\end{document}
